\documentclass[../Main.tex]{subfiles}

\begin{document}

\section{Reflexion} \label{sec:Reflexion}

\subsection{Reflexion des Arbeitsprozesses}

Nach rund einen ganzen Jahr von Konzipieren, Recherchieren und Programmieren können wir nun die Maturaarbeit Project Polearm beenden. Das Jahr war gefüllt mit ... (Momenten???). Was haben wir nun gemacht, und was haben wir daraus gelernt? Während wir die Einträge für unsere Arbeitsjournale verfasst haben, haben wir diese Frage auf einer kleineren, täglichen Ebene beantwortet, und nun wird diese Reflektion dies nun in hinbezug auf das ganze Projekt erschliessen.

Die Maturaarbeit haben wir in rund vier Phasen unterteilt, nähmlich die Planung, Recherche, Entwicklung und zuletzt der Playtest inklusive Evaluierung. Manche von denen Phasen sind parallel zueinander gelaufen, während andere voneinander abhängig gewesen sind. Was haben die Phasen für sich aber beansprucht?

Fangen wir mit der Planung an, denn aus der Planung stammt ja alles. Dafür müssen wir etwas zurückkehren, noch vor den Start der Maturaarbeit. Die Idee, einen Spiel zu entwickeln, ist uns relativ früh in den Sinn gekommen, denn der Projektunterricht in der dritten Klasse und die Maturaarbeit haben wir als eine Chance gesehen, etwas Grosses machen zu können(, ein richtiges Produkt zu entwickeln). Im Projektunterricht haben wir somit angefangen, Recherche über die Thematik des Spiels zu machen. Und als nun die Maturaarbeit kam, hatten wir schon eine Säule des Spiels fertig. Wir fühlten uns bereit, mit den anderen fortzufahren.

Wir sind also mit einer festen Vision in die Maturaarbeit hineingestiegen, und haben früh angefangen, Spielsystemen zusammen zu diskutieren und zu (prototypen?). Etwas zu früh vielleicht, da wir noch nicht wussten, dass wir eine Erkenntnisfrage brauchten. Als wir dies mitbekamen, mussten wir unsere Prioritäten umstellen, und damit auch unsere Erwartungen. Nun mussten wir überlegten, wie wir mit unseren geplanten Mechaniken die Erkenntnisfrage beantworten wollen.  Da wir auch unsere Erkenntnisfrage ordentlich und belegt beantworten wollen, mussten wir uns auch mit der Recherche auseinandersetzen. (Erkenntnisfrage erläutern irgendwo?)

Nach der Konzeption blieben uns bis zur Zwischenpräsentation nun noch vier Monate. Wir hatten eine grosse Freiheit verspürt, denn einen groben Plan für das Spiel hatten wir schon. Es gab eine grosses Buffet von Sachen, die wir machen konnten, die Auswahl stand zwischen etwas technisches oder literarisches. Die Fallhöhe war zudem gering: Was auch immer wir machten, nichts musste verworfen werden, das Projekt war noch ein junger und wild wachsendes Setzling, und alles kann man schliesslich einmal ausprobieren. Zu diesen Punkt hingen die Arbeitsschritte noch kaum voneinander ab. Rückblickend müssen wir zurückgeben, dass wir nicht ganz mit den Vertrag im Auge gearbeitet haben, denn wir waren sicher, dass die Vertragsziele ziemlich erreichbar waren. Dass diese naive Selbstsicherheit und unabhängige Arbeiten aufging, müssen wir zum teil den Glück bedanken, aber ist die Natur eines Projektanfangs nicht so? Man hört oft von Projektretroperspektiven, dass Projekte am Anfang immer etwas wilder sind, vieleicht weil man Anfang noch eine blinde Motivation hatte, denn man weiss nie wie viel Arbeit es am Ende sein. Könnten wir noch einmal beginnen, würden wir es genauso angehen, mit dem einzigen Unterschied, dass wir wüssten, dass wir diese rapide, motivierte Anfangsphase eine Frist hat.

Bis zur Zwischenpräsentation spielte die Erkenntnisfrage eine untergeordnete Rolle, denn die dafür relevanten Systeme beruhten auf einer technischen Basis, die wir noch nicht ganz hatten: einem Spielercharakter mit Animationen und Steuerung, sowie die Anfänge eines Gegners. Unsere "Freiheit" in der frühen Planungs- und Entwicklunsphase hat sich durch die Implementierung von Spielexternen Systeme wie Menus, Einstellungen und weiteres ausgedrückt. Unser Spielercharakter war auch weitaus Fortgeschrittener, als der Vertrag verlangt hätte. Dies ist uns in der Bewertung auch zu einen Krikpunkt geworden: Nicht weiter ausführen, sondern jetzt spezifizieren. Zum einen machte uns dieser Teil am meisten Freude, zum anderen mussten wir uns zuerst eine Vorstellung von unserem eigenen Spiel bilden. Die Recherche hat später angefangen als die Entwicklung, und war dementsprechend auch in einer (unterentwickelten?) Zustand. Das war aber rückblickend die richtige Entscheidung. Erst nachdem wir die ersten Systeme gebaut hatten, wussten wir, worauf wir in anderen Spielen achten mussten, und konnten erst recht eine gute Recherche machen. Bei einem nächsten Projekt würden wir die Reihenfolge davon abhängig machen, wie konkret unsere Vorstellungen sind und wie gross unser Kenntnisstand ist: Haben wir beides, recherchieren wir zuerst. Fehlt eines davon, sammeln wir zuerst eigene Erfahrungen, also genau das, was wir diesmal eher unbeabsichtigt getan haben.

Weitere Kritikpunkte der Zwischenpräsentation betrafen auch die begleitenden Texte, die wir ab den Punkt schon hatten: Er war nicht klar, ob es sich bei unsere Aussagen um eigene Einschätzungen oder um belegte Aussagen handeln, dies müssten wir für die finale Arbeit umformulieren. Unsere KI-Nutzung war zwar dokumentiert, aber nicht so, wie man es erwartet hätte. Beim Umsetzen hatten wir zunächst Schwierigkeiten, denn hätten wir die bemängelten Punkte verinnerlicht gehabt, hätten wir sie von Anfang an richtig gemacht. Deshalb stellten wir unserer Betreuungsperson so lange Fragen, bis wir glaubten, diese Unklarheiten aufgeklärt haben.

Und ab diesen Punkt haben sich die ersten Probleme hineingeschlichen. Wir haben unsere Featureliste betrachtet, und manche Funktionalitäten wir zwar wollten, die für die Erkenntnisfrage aber wenig relevant waren, etwa ein Stylemeter wie in \game{Devil May Cry}\autocite{capcom_devilmaycry_2001} oder eine zweite, auswechselbare Waffe, mussten wir schon jetzt streichen oder stark eingrenzen, und wir haben realisiert, dass wir priorisieren müssen, was für die Arbeit essenziell ist. Dazu kam ein begründet geplanter Playtest, und die Checkliste aus dem Vertrag hatte noch viele offene Punkte. Das Fertigstellen der vorgesehenen Vertragspunkte waren jedoch nicht das Problem, die Qualitätssicherung hat uns den grössten Strich durch die Rechnung gemacht, weil wir möglichst wenige Kompromisse eingehen wollten. Im Nachhinein wäre es interessant gewesen zu erfahren, wie 80 oder 85 statt 100 Prozent unseres Anspruchs bewertet worden wären (formuliere das besser um). Da die Zwischenbewertung nicht in die Note der Maturitätsarbeit einfliesst, ist das eine verpasste Chance. Andererseits konnten wir Erfahrung mit etwas leichten Crunchen sammeln, die uns am Ende ermöglicht hat, mit genügend Zeit und 95 von 100 Prozent unserem Potentials herauszuholen.

Nach der Zwischenpräsentation hatten wir wieder genug Zeit, und die Abgabe war weit genug entfernt. Wir haben unseren Blick auf etwas weitere Horizonte geworfen, aber gleichzeitig nahm das anfängliche Tempo ab, mit dem wir Neues produzierten. Zum einen war das Projekt inzwischen so gross, dass Neuerungen weniger Unterschied machten, sie fühlten sich weniger an. Zum anderen war unser anfänglicher Eifer zwar weit davon entfernt, erloschen zu sein, doch das Projekt zu pflegen und Aufgaben zu erledigen, die uns weniger Spass machten und sich mittlerweile angehäuft hatten, setzte uns langsam zu. Wir nahmen deshalb beide etwas Abstand vom Projekt und verfolgten andere Interessen, was möglich war, weil wir gut vorgearbeitet hatten. Danach kehrten wir gestärkt zurück, zumal wir die nächste Deadline im Blick hatten und uns ohnehin keine andere Wahl blieb. (Willst du hier einen satz schreiben? Wenn du in einen von deinen side quests gemerkt hast, dass du in etwas besser geworden bist, schreibe es hin, und sag mir bescheid, damit ich auch ein beispiel geben kann)

Mittlerweile hat die Recherche richtig angefangen, und unsere Rollen im Projekt haben sich konsolidiert (referenz zum späteren kapitelabschnitt?). Die einzelnen Arbeitsschritte hingen stark voneinander ab, und wir konnten in unserer Partnerarbeit nicht mehr einfach tun und lassen, was wir wollten, wir brauchten einen anderen Arbeitsprozess. Wir arbeiteten deshalb an möglichst grossen, voneinander unabhängigen Blöcken und vereinbarten ein Datum, wenn sich unsere Arbeiten kreuzten. In dieser Zeit entstand auch fast die gesamte restliche Recherche. Dieser Prozess funktionierte sehr gut, wenn beide rechtzeitig fertig wurden, was allerdings nicht immer der Fall war. (Künftig würden wir deshalb einen dritten, grösseren Arbeitsstrang offenhalten, an dem weiterarbeiten kann, wer gerade durch die andere Person blockiert ist). (Die aussage ist riskant und "weniger" wert, als wenn man larpt und sagt, ja wir werden versuchen leicht unabhäniger zu arbeiten, und mehr tatsächlich zusammen, WÄHREND noch kleinere seitenaufgaben frei sind, wenn die andere person blockiert ist usw.)

Nach Abschluss der Recherche wussten wir nun genau, welche Systeme noch implementiert werden mussten, und mit der näher rückenden Abgabe hatten wir nun auch die Beantwortung der Erkenntnisfrage im Blick. Hier beginnt der wirklich gut geplante Teil unserer Arbeit. Da wir nun ein Fundament hatten, das in eine klare Richtung ging, konnten wir zielgerichtet auf die Erkenntnisfrage hinarbeiten. Dies war jedoch nicht immer positiv, da wir unsere Anfängliche Freiheit nun verloren haben. Manches der Passion am Anfang wurde nun Fleissarbeit. Trotzdem müssen wir zugeben, dass dieser Teil ist uns besonders gelungen, weil wir auf Hinweis unserer Betreuungsperson einen Plan bis zur Abgabe erstellt haben (wiederholung, formuliere mit chatty um). Dass dieser Plan gut wurde, verdanken wir den Erfahrungen aus der Zwischenpräsentation, bei der wir bereits gelernt hatten, wie man auf eine Abgabe hinarbeitet und plant.

Unseren ersten fertigen Export hatten wir am 25.~Oktober. Technisch war das Spiel damit nahezu fertig, doch es hat noch die Polierung gefehlt. Wir konnten bereits die Früchte unserer Arbeit und das Potenzial sehen trotz seiner rauen Ecken. Die Mechaniken, um unsere Rechercheergebnisse umzusetzen, waren schon vorhanden, doch konsequent und beobachtbar implementiert waren sie noch nicht. Nun mussten wir an ein den Stellschrauben nachjustieren und die fehlenden Aspekte ergänzen, um das Spiel wurde unserer Recherche gerecht zu machen und für die Playtests bereit zu machen. Diese Erfahrung hat uns gezeigt, dass der Teufel im Detail steckt. Bis dahin hatten wir Recherche und Prototyp getrennt betrachtet, aber die Schnittstelle, die zwischen den Beiden nun entstehen musste, war eine Detailarbeit. Die unabhängige Arbeit am Anfang fädelten sich nun langsam zusammen zu der selben Schnur, und die Arbeit wurde vereint. (Ziehe hier eine neue Conclusion)

Gegen Ende der Arbeit führten wir die Playtests durch. Sie verliefen unserer Meinung nach gut, da wir bei der Analyse von \game{NieR: Automata} und \game{Ninja Gaiden} bereits selbst Spieletests durchgeführt hatten, wir wussten, was man in einen Playtest erwarten könnte und haben unsere Planung dementsprechend angepasst. In der Schlussphase bereiteten uns eher Formalitäten Sorgen. Uns war bewusst, dass es knapp wird, wenn wir uns nicht an den Plan halten, auch wenn es noch nach genug Zeit aussah. Daher haben wir Pufferzeiten eingeplant, die nun uns nützlich wurden. Es war genau diese gezielte Planung, die uns zwar etwas Passion am Projekt geraubt hat, aber letztenendes uns vieles Stress entzogen hatte. Dieser sind wir dankbar, dass diese Phase verlief unserer Einschätzung gut lief. In Zukunft wäre es Interessant zu erkunden, wie man Passion und Disziplin miteinander von Anfang an Balancieren kann.

Zusammengefasst, der Entwicklungsprozess von Project Polearm war einer gefüllt mit Motivation und wilde Passion, die wir mit der ruhigen und geziehlten Planung balanciert haben. Am Anfang durften wir erleben, wie ein Projekt wie ein Gebüsch wächst, der zum Schluss zu einen gerichteten und langsamen wachsenden Baum wurde. (etwas yapping hier maybe hast du ne idee?) Aber was ist mit unseren Anfänglichen Ambitionen? Hier sagen wir, dass obwohl diese nicht so erfüllt worden waren, wie wir vor einen Jahr erwartet hätten, doch stattdessen haben wir etwas viel wichtigers gesammelt. So wie das Projekt an Reife und Präzision gewonnen hat, so haben wir dies auch getan. Wir haben gelernt, wie man über Programme strukturell nachdenkt und wie man Systeme zusammenführt, und ebenfalls wie man ein grosses Projekt zusammenführt und eine gemeinsame Teamarbeit koordiniert. Wir haben beobachtet, wie unterschiedliche Formen des Projekt-Managements zu drastisch anderen Arbeitsprozessen führen kann. Auch wenn wir es am Anfang vieleicht noch verspöttet hätten, unsere Reise ist zu unseren Ziel geworden, ein Ziel, dessen Lehren wir in usneren Leben mit uns tragen werden.

\subsection{Aufgabenteilung und Eigenleistung} \label{sec:Aufgabenteilung}

In dieser Arbeit haben wir unsere Aufgaben aufgeteilt. Manches haben wir zusammen gemacht, manches ganz allein. In Tabelle~\ref{tab:Aufgabenteilung} ist grob festgehalten, was wir hauptsächlich gemacht haben bzw.\ wo der grösste Anteil unserer Arbeit hineingeflossen ist. Genaue Angaben finden sich in unseren Arbeitsjournalen.

\begin{table}[H]
    \centering
    \small
    \renewcommand{\arraystretch}{1.3}
    \begin{tabular}{|>{\raggedright\arraybackslash}p{2.6cm}|>{\raggedright\arraybackslash}p{5.8cm}|>{\raggedright\arraybackslash}p{5.8cm}|}
        \hline
        & \textbf{Rodin Arslan} & \textbf{Zhile Zhu} \\
        \hline
        \textbf{Prototyp} &
        \begin{minipage}[t]{\linewidth}\raggedright
            \begin{itemize}[nosep, leftmargin=*, label=--]
                \item Gamedesign: Anforderungsdefinition, Verfeinerung der Feature-Konzeption sowie der Umsetzung, Fine-Tuning des Spielgefühls
                \item Technische Umsetzung von Menüs, Einstellungen und Shadern
                \item Assets: Beschaffung von 3D-Modellen, Animationen und Texturen, Erstellung des Levels und dessen Modells, Integration ins Spiel
            \end{itemize}\vspace{0.4em}
        \end{minipage}
        &
        \begin{minipage}[t]{\linewidth}\raggedright
            \begin{itemize}[nosep, leftmargin=*, label=--]
                \item Gamedesign: Feature-Konzeption
                \item Technische Umsetzung der Features sowie Sauberkeit und Einheitlichkeit des Programms
                \item Assets: Beschaffung von Musik und SFX, Erstellung von Spieler- und Gegner-Aktionsdaten, Integration ins Spiel
            \end{itemize}\vspace{0.4em}
        \end{minipage}
        \\
        \hline
        \textbf{Recherche und Analyse} &
        \begin{minipage}[t]{\linewidth}\raggedright
            \begin{itemize}[nosep, leftmargin=*, label=--]
                \item Leveldesign-Recherche
                \item Referenzspielanalyse (\game{NieR: Automata})
            \end{itemize}\vspace{0.4em}
        \end{minipage}
        &
        \begin{minipage}[t]{\linewidth}\raggedright
            \begin{itemize}[nosep, leftmargin=*, label=--]
                \item Referenzspielanalyse (\game{Ninja Gaiden Sigma Plus})
            \end{itemize}\vspace{0.4em}
        \end{minipage}
        \\
        \hline
        \textbf{Playtests} &
        \begin{minipage}[t]{\linewidth}\raggedright
            \begin{itemize}[nosep, leftmargin=*, label=--]
                \item Planung, Durchführung und Auswertung
            \end{itemize}\vspace{0.4em}
        \end{minipage}
        & \\
        \hline
        \textbf{Dokumentation} &
        \begin{minipage}[t]{\linewidth}\raggedright
            \begin{itemize}[nosep, leftmargin=*, label=--]
                \item Begleitbericht
            \end{itemize}\vspace{0.4em}
        \end{minipage}
        &
        \begin{minipage}[t]{\linewidth}\raggedright
            \begin{itemize}[nosep, leftmargin=*, label=--]
                \item Technische Dokumentation
            \end{itemize}\vspace{0.4em}
        \end{minipage}
        \\
        \hline
    \end{tabular}
    \caption{Aufgabenteilung: Schwerpunkte der beiden Gruppenmitglieder}
    \label{tab:Aufgabenteilung}
\end{table}

\subsection{Verwendung von KI} \label{sec:ReflexionKI}

\game{Project Polearm} war ein grosses Informatikprojekt, da muss doch viel KI im Spiel gewesen sein, oder? Eigentlich nicht. Unsere Philosophie war, dass KI nichts erschaffen soll, sondern nur ein Werkzeug ist, das besonders gut darin ist, Informationen aufzubereiten. Natürlich ist KI auch gut darin, Code zu schreiben und zu bereinigen. Wir halten unseren Umgang mit KI für gesund, weil er ihre Stärken ausnutzt, wie unser Vorgehen zeigt.

Konkret haben wir die KI beim Brainstorming und bei der Planung nach ihrer Einschätzung gefragt. Für die Recherche haben wir selbst über Google Scholar gesucht, unsere Bücher hat uns allerdings die KI vorgeschlagen, als wir nach Material gefragt haben. Im Code haben wir KI nirgends eingesetzt, wo es für die Erkenntnisfrage relevant war, und eher für Prototypen als für fertige Systeme. Am fertigen Spiel ist KI nur bei den Shadern beteiligt, die sie für uns portiert oder selbst geschrieben hat, sowie beim Soundsystem. Beim Arbeitsjournal wurde KI genutzt, um unsere Einträge so zu strukturieren, wie es unsere Betreuungsperson verlangt hatte.

Im Begleitbericht hat die KI die Rückverweise gesetzt und ein Glossar für Begriffe erstellt, die unverständlich sein könnten. Neben dem Glossar wurden einige wenige weitere Texte mit KI formuliert. Diese kann man in der KI-Deklaration (Abschnitt~\ref{sec:KIDeklaration}) nachschlagen. Zudem hat sie den gesamten Begleitbericht auf Rechtschreibung geprüft und nach logischen Fehlern gesucht, die wir über rund 100 Seiten hinweg selbst nicht im Auge behalten konnten. Die Playtest-Protokolle hat sie anonymisiert und geschlechtsneutral gemacht, und die \LaTeX-Formatierung haben wir ihr diktiert, worauf sie diese für uns umgesetzt hat.

Die KI hatte aber auch ihre Grenzen. Über die Dauer des Projekts haben wir mehrere Chatverläufe verloren, sodass der KI manchmal der Kontext fehlte und wir ihn in einem neuen Chat erneut aufbauen mussten. Zudem kann die KI weder wissen, was in unseren Köpfen vorgeht, noch den Zusammenhang des gesamten Projekts kennen, den wir kennen.

Als Fazit sehen wir die grösste Gefahr bei der Nutzung von KI darin, das Projekt an sie zu verlieren, weil der Eigenanteil zu klein und der Anteil der KI zu gross wird. Wir haben unsere Nutzung von KI zudem eingeschränkt, da wir unsere Kreativität bewahren wollten und es uns wichtig war, inhaltliche Entscheidungen selbst zu treffen. Inhaltlich war KI für uns nur ein Mittel, unsere Kreativität anzuregen.

\end{document}
